% =====================================================================
% tesi/capitoli/A3-appendice-algebra.tex
% =====================================================================
% copre: docs/40-appendice-matematica.md#3-1-fattoriale
% copre: docs/40-appendice-matematica.md#3-2-permutazione
% copre: docs/40-appendice-matematica.md#3-3-gruppo-abeliano-e-la-commutatività-da-cui-segue-l-invarianza-del-checksum
% copre: docs/40-appendice-matematica.md#3-4-anello-dei-resti-modulo-n
% copre: docs/40-appendice-matematica.md#3-5-aritmetica-modulo-due-e-lo-xor-come-somma
% copre: docs/40-appendice-matematica.md#3-6-teorema-cinese-del-resto
% copre: docs/40-appendice-matematica.md#3-7-coefficiente-binomiale
% copre: docs/40-appendice-matematica.md#3-8-combinazioni-con-ripetizione
% copre: docs/40-appendice-matematica.md#3-9-principio-di-inclusione-ed-esclusione
% copre: docs/40-appendice-matematica.md#3-10-principio-dei-cassetti
% copre: docs/40-appendice-matematica.md#3-11-sistema-numerico-fattoriale-e-codice-di-lehmer
% verificato-al-commit: 3439cc8
% ---------------------------------------------------------------------

\chapter{Nozioni di algebra e di combinatoria}
\label{app:algebra}

Questa appendice raccoglie le nozioni che sostengono i due risultati più sorprendenti del capitolo~\ref{cap:analisi}, cioè l'invarianza del checksum per permutazione e l'identificazione della tabella delle ventiquattro permutazioni con il codice di Lehmer, oltre agli strumenti di conteggio impiegati nell'appendice~\ref{app:geometria}.

\section{Fattoriale}

Il fattoriale di un intero non negativo $n$, scritto $n!$, è il prodotto degli interi da uno a $n$, con la convenzione $0! = 1$. Conta le disposizioni in fila di $n$ oggetti distinti.

La nozione esiste perché il fattoriale è la cardinalità di base di ogni conteggio di disposizioni, e perché la sua crescita, più rapida di qualunque esponenziale, è la ragione per cui certi spazi si esplorano esaustivamente e altri no.

L'esempio da svolgere è quello che serve al lavoro: $4! = 24$, cioè il numero di ordinamenti delle quattro sottostrutture di un esemplare della terza generazione, e dunque il numero di righe della tabella di permutazione del capitolo~\ref{cap:cifratura}. La convenzione $0! = 1$ non è arbitraria: è ciò che rende valida la formula del coefficiente binomiale nei casi estremi e ciò che chiude la decodifica del codice di Lehmer nell'ultima sezione di questa appendice.

\section{Permutazione}

Una permutazione di un insieme finito è una funzione biiettiva dell'insieme in sé, cioè un riordinamento che non perde né duplica elementi. Le permutazioni di $n$ elementi sono $n!$ e si compongono fra loro come funzioni.

La nozione esiste perché distingue una trasformazione che riorganizza da una che altera. Una permutazione è invertibile, dunque non distrugge informazione, e questo fatto è impiegato nel capitolo~\ref{cap:analisi} in due direzioni opposte: la permutazione delle sottostrutture non aggiunge sicurezza perché non aggiunge incertezza, e la permutazione dei blocchi non viene rilevata dal checksum perché la somma non dipende dall'ordine.

L'esempio da svolgere è la permutazione delle quattro sottostrutture, che il gioco determina dal resto della divisione del valore di personalità per ventiquattro. Poiché la funzione è biiettiva, chi conosce quel resto ricostruisce l'ordine originario applicando la permutazione inversa, e la procedura di lettura risulta simmetrica a quella di scrittura: è la proprietà su cui poggia la verifica di simmetria descritta nel capitolo~\ref{cap:collaudo}.

\section{Gruppo abeliano, e l'origine dell'invarianza del checksum}

Un gruppo è un insieme dotato di un'operazione associativa con elemento neutro, in cui ogni elemento ammette inverso. Il gruppo si dice abeliano, o commutativo, quando l'ordine degli operandi è indifferente, cioè quando $a + b = b + a$ per ogni coppia di elementi.

La nozione esiste, in questo lavoro, perché nomina esattamente la proprietà responsabile di un difetto reale, e nominarla converte un'osservazione empirica in una dimostrazione. Il checksum della terza generazione somma parole nel gruppo abeliano degli interi modulo $2^{16}$; la somma di un insieme di addendi in un gruppo abeliano non dipende dall'ordine in cui gli addendi sono presi; dunque una permutazione dei blocchi produce la medesima somma. Non con probabilità elevata: sempre.

L'esempio da svolgere è minimo e verificabile a mente, giacché la somma di $\mathtt{0x1234}$ e $\mathtt{0x5678}$ vale $\mathtt{0x68AC}$ in entrambi gli ordini. Il caso reale è la tabella delle ventiquattro permutazioni, dove le settecentodiciannove alterazioni consistenti in un riordinamento delle sottostrutture superano il controllo con probabilità unitaria, contro la probabilità $2^{-16}$ attribuibile a un'alterazione generica. È la ragione per cui il capitolo~\ref{cap:collaudo} conclude che il confronto con un'implementazione indipendente non è ridondante rispetto alla verifica di simmetria: le due difese sono cieche al medesimo genere di errore.

\section{Anello dei resti modulo $n$}

L'insieme $\{0, 1, \dots, n-1\}$ con addizione e moltiplicazione definite prendendo il resto del risultato ordinario forma l'anello dei resti modulo $n$. Ogni operazione resta all'interno dell'insieme, che è precisamente il comportamento di un registro di larghezza fissa quando trabocca.

La nozione esiste perché il traboccamento non è un errore da evitare ma un'operazione ben definita di cui conoscere le regole: un'addizione a \SI{16}{bit} che trabocca non produce un valore indefinito, produce la somma modulo $2^{16}$, e i checksum dei giochi si appoggiano a questo fatto anziché combatterlo.

L'esempio da svolgere è il checksum della terza generazione, che somma parole di \SI{16}{bit} conservandone i soli sedici bit inferiori: poiché $65535 + 3 = 65538$ e $65538 \bmod 65536 = 2$, la somma di $\mathtt{0xFFFF}$ e $\mathtt{0x0003}$ vale $\mathtt{0x0002}$. Il confronto con la prima generazione è istruttivo: là il controllo è a \SI{8}{bit} e la somma è nell'anello modulo $2^8$, dunque i valori distinti che il controllo può assumere sono $256$ anziché $65\,536$, ed è questa la ragione per cui la sua capacità di rilevazione è inferiore di due ordini di grandezza, come mostra il capitolo~\ref{cap:integrita}.

\section{Aritmetica modulo due, e lo XOR come somma}

L'aritmetica modulo due ha due soli elementi e in essa l'addizione coincide con l'or esclusivo, poiché $1 + 1 = 0$. Ogni elemento è il proprio opposto, dunque sommare due volte la medesima quantità la annulla.

La nozione esiste perché spiega in una riga tutte le proprietà del cifrario a chiave scorrevole e della mascheratura delle quantità: cifratura e decifratura sono la medesima operazione, e ciò non è una comodità implementativa ma una conseguenza algebrica.

L'esempio da svolgere è la mascheratura dello zaino descritta nel capitolo~\ref{cap:smeraldo}. La quantità in chiaro $q$ è memorizzata come $q \oplus k$ con $k$ la chiave di sicurezza, e in lettura si applica di nuovo la chiave ottenendo $(q \oplus k) \oplus k = q \oplus (k \oplus k) = q$. Dalla medesima proprietà discende il difetto misurato nel capitolo~\ref{cap:analisi}: se due dati distinti sono mascherati con la medesima chiave, il loro or esclusivo elide la chiave e lascia $d_1 \oplus d_2$, cioè una relazione fra i chiari ottenibile senza conoscere $k$.

\section{Teorema cinese del resto}

Il teorema stabilisce che, noti i resti di un intero rispetto a moduli a due a due coprimi, l'intero è determinato in modo unico modulo il prodotto dei moduli. La forma rilevante qui è la sua conseguenza negativa: se i moduli non sono coprimi, la determinazione non è unica e l'informazione dei due resti si sovrappone.

La nozione esiste perché stabilisce quando due vincoli su resti diversi sono indipendenti e quando invece si intralciano, che è la domanda posta da chi voglia costruire un valore di personalità soddisfacente più condizioni simultaneamente.

L'esempio da svolgere è quello del lavoro. Il sesso dipende dal resto modulo $256$, la natura dal resto modulo $25$, la permutazione dal resto modulo $24$. I moduli $256$ e $25$ sono coprimi, essendo potenze di due e di cinque, dunque i due vincoli sono indipendenti ed esiste un valore che li soddisfa entrambi, unico modulo $6400$. I moduli $256$ e $24$ condividono il fattore otto, dunque quei due vincoli non sono indipendenti e alcune combinazioni sono irrealizzabili. È la ragione per cui il risolutore di vincoli discusso nel capitolo~\ref{cap:conversione} non può trattare i vincoli come una lista da soddisfare uno alla volta.

\section{Coefficiente binomiale}

Il coefficiente binomiale conta i sottoinsiemi di $k$ elementi estraibili da un insieme di $n$ elementi distinti, senza riguardo all'ordine:

\begin{equation}
\binom{n}{k} = \frac{n!}{k!\,(n-k)!} .
\end{equation}

La nozione esiste perché separa il conteggio delle scelte da quello delle disposizioni: quando l'ordine è irrilevante, contare le disposizioni sovrastima di un fattore $k!$, e il coefficiente binomiale è la correzione di quel fattore.

L'esempio da svolgere è un caso minimo verificabile a mano, $\binom{4}{2} = 6$, cioè i sei modi di scegliere due sottostrutture su quattro. Nel lavoro il coefficiente compare nella distribuzione binomiale dell'appendice~\ref{app:probabilita} e nel conteggio dei punti interi dell'appendice~\ref{app:geometria}.

\section{Combinazioni con ripetizione}

Il numero di modi in cui $s$ unità indistinguibili si distribuiscono fra $d$ contenitori distinti, ammettendo contenitori vuoti, vale

\begin{equation}
N(s,d) = \binom{s + d - 1}{d - 1}.
\end{equation}

La nozione esiste perché conta le soluzioni intere non negative di un'equazione a somma fissata, cioè le configurazioni di ogni sistema in cui una risorsa totale si ripartisce fra più campi. La dimostrazione classica è l'argomento delle barre e delle stelle: si dispongono le $s$ unità in fila e si inseriscono $d-1$ separatori, e ogni disposizione dei separatori fra i $s + d - 1$ posti individua una ripartizione.

L'esempio da svolgere è la forma impiegata nel lavoro. La somma dei sei valori di allenamento di un esemplare della terza generazione non supera $510$ e ciascuno non supera $252$: il conteggio delle configurazioni ammissibili parte da $N(s,6)$ per ogni somma $s$ e sottrae le violazioni del tetto individuale con il metodo dell'appendice~\ref{app:geometria}. Il risultato, $22\,858\,382\,491\,812$ punti, corrisponde a \SI{44,38}{bit} e costituisce la capacità informativa dell'insieme di arrivo della conversione.

\section{Principio di inclusione ed esclusione}

Il principio calcola la cardinalità di un'unione sommando le cardinalità dei singoli insiemi, sottraendo quelle delle intersezioni a due, aggiungendo quelle a tre, e proseguendo con segni alterni. Nella forma impiegata qui viene usato al contrario, per contare gli elementi che non violano alcun vincolo, sottraendo dal totale le violazioni con la medesima alternanza.

La nozione esiste perché sommare le violazioni le conta più volte quando due violazioni possono coesistere, e il principio è la contabilità esatta di quella sovrapposizione: omettere il termine di riaggiunta produce un risultato inferiore al vero, ed è l'errore tipico.

L'esempio da svolgere è il conteggio dell'allenamento. Le configurazioni con somma $s$ sono $N(s,6)$; quelle in cui un campo determinato supera $252$ si contano collocandovi $253$ unità e ripartendo le rimanenti, cioè $N(s-253,6)$, e i campi sono sei; le violazioni doppie richiedono $506$ unità, dunque compaiono soltanto per $s \ge 506$, e sono $\binom{6}{2} = 15$ termini $N(s-506,6)$; le violazioni triple richiederebbero $759$ unità e non esistono per $s \le 510$. La somma alternata si arresta quindi al secondo termine, e questa chiusura anticipata è ciò che rende il risultato esatto anziché approssimato.

\section{Principio dei cassetti}

Il principio afferma che collocando $n$ oggetti in $m$ contenitori con $n > m$ almeno un contenitore riceve più di un oggetto e, nella forma quantitativa, che almeno un contenitore riceve almeno $\lceil n/m \rceil$ oggetti.

La nozione esiste perché costituisce la dimostrazione di impossibilità più economica disponibile: non richiede di esaminare la funzione, ma soltanto di contare gli insiemi. Ed è la forma di tutte le dimostrazioni di perdita di informazione di questo lavoro.

L'esempio da svolgere è la conversione dell'allenamento. Le configurazioni di partenza sono $2^{80}$, i punti di arrivo meno di $2^{45}$: qualunque funzione fra i due insiemi manda almeno $2^{35}$ configurazioni distinte nel medesimo punto, dunque nessuna può essere iniettiva e la fedeltà è impossibile. La conclusione è indipendente dalla formula, e questo è il punto: il principio dimostra un limite del formato e non un difetto dell'implementazione.

\section{Sistema numerico fattoriale e codice di Lehmer}

Nel sistema numerico fattoriale la cifra di posizione $i$, contata da destra a partire da uno, ha peso $i!$ e assume valori da zero a $i$. Ogni intero non negativo minore di $n!$ ammette una e una sola rappresentazione con $n-1$ cifre. Il codice di Lehmer è la corrispondenza che ne discende fra gli interi da $0$ a $n!-1$ e le permutazioni di $n$ elementi: si legge la prima cifra come indice dell'elemento da estrarre dalla lista ordinata, si rimuove quell'elemento e si procede con la cifra successiva sulla lista residua.

La nozione esiste perché fornisce una numerazione canonica delle permutazioni, cioè un modo di passare da un indice a una permutazione senza tabelle. Riconoscere che una tabella incorpora quella numerazione consente di sostituirla con un algoritmo e, prima ancora, di verificarla anziché fidarsene.

L'esempio da svolgere è la verifica compiuta in questo lavoro. Per $n = 4$ gli indici vanno da $0$ a $23$; si consideri l'indice $7$. La divisione per $3! = 6$ dà quoziente $1$ e resto $1$, dunque si estrae il secondo elemento della lista ordinata, cioè $B$. Il resto $1$ diviso $2! = 2$ dà quoziente $0$ e resto $1$, dunque si estrae il primo elemento residuo, cioè $A$. Il resto $1$ diviso $1! = 1$ dà quoziente $1$, dunque si estrae il secondo elemento residuo, cioè $D$; resta $C$. La permutazione è $B\,A\,D\,C$, che coincide con la riga di indice $7$ della tabella del gioco. La verifica è stata condotta su tutti i ventiquattro indici, e l'unicità della rappresentazione si dimostra per conteggio: le rappresentazioni possibili sono $1 \cdot 2 \cdot 3 \cdots n = n!$, tante quante le permutazioni, e la corrispondenza è iniettiva per costruzione, dunque biiettiva.
